% =============================================================================
% Capítulo 3 — Metodologia
% =============================================================================
\chapter{Metodologia}
\label{cap:metodologia}

Este capítulo descreve as escolhas metodológicas que orientaram o
desenvolvimento do presente trabalho. São apresentados a concepção do
protótipo jogável, o ambiente técnico e as ferramentas utilizadas, o
pipeline completo de coleta, extração e treinamento, o protocolo de avaliação
quantitativa e, por fim, as decisões de escopo que delimitaram a contribuição
do trabalho.

% -----------------------------------------------------------------------------
\section{Concepção do Protótipo (\textit{Vertical Slice})}
\label{sec:vertical-slice}

A estratégia de desenvolvimento adotada foi a construção de um
\textit{vertical slice}: um recorte funcional do jogo que implementa, de
ponta a ponta, todas as camadas necessárias para validar a proposta --- desde
a jogabilidade até a exportação de métricas --- sem a pretensão de representar
um produto comercial completo. Essa abordagem é recorrente na indústria de
jogos quando o objetivo é demonstrar a viabilidade de um conceito em escopo
controlado \cite{yannakakis2018artificial}.

O protótipo, denominado \textit{Chronicles Of The Lost Word}, consiste em um
jogo 2D de ação construído no motor Unity. O recorte implementado concentra-se
exclusivamente em uma arena de combate contra um chefe (\textit{boss}),
excluindo elementos como exploração de fases, sistemas de progressão,
narrativa interativa ou interfaces de menu. Essa restrição é intencional: o
escopo reduzido permite isolar a variável de interesse --- o comportamento do
chefe --- e avaliar o impacto do pipeline de proveniência e \textit{Reward
Shaping} sem interferência de sistemas periféricos.

\subsection{O Personagem Jogável (Elian)}
\label{subsec:elian}

O personagem controlado pelo jogador, Elian, possui um conjunto de mecânicas
de ação que define o espaço de possibilidades da interação com o chefe:

\begin{itemize}
    \item \textbf{Movimento horizontal}: deslocamento à esquerda e à direita
          no plano da arena;
    \item \textbf{Pulo} (\textit{PlayerJump}): salto vertical com detecção
          de solo via \textit{ground check};
    \item \textbf{Ataque} (\textit{PlayerAttack}): golpe corpo a corpo com
          área de efeito definida por \texttt{OverlapBoxAll}, aplicando dano
          em alvos que implementam a interface \texttt{IDamageable};
    \item \textbf{Dash} (\textit{PlayerDash}): avanço rápido horizontal com
          breve período de invulnerabilidade;
    \item \textbf{Vida} (\textit{PlayerHealth}): sistema de pontos de vida
          com dano recebido, morte e eventos associados.
\end{itemize}

Essas três ações ofensivas e evasivas --- pulo, ataque e \textit{dash} ---
constituem o vocabulário de ações registrado pelo sistema de proveniência e,
posteriormente, utilizado na extração de sequências vitoriosas. A escolha
desse conjunto restrito foi deliberada: um vocabulário de ação compacto
facilita a identificação de padrões causais nos grafos de proveniência e
mantém o espaço de sequências possíveis tratável computacionalmente.

\subsection{A Arena de Combate}
\label{subsec:arena}

A arena consiste em um espaço horizontal delimitado por paredes invisíveis,
com plataforma plana e sem obstáculos adicionais. Esse \textit{design}
minimalista garante que as diferenças de comportamento entre os chefes FSM e
PPO reflitam exclusivamente suas estratégias de combate, sem influência de
variáveis ambientais como plataformas elevadas, armadilhas ou itens coletáveis.

Cada sessão de combate é gerenciada pelo componente \texttt{ArenaManager}, que
controla o ciclo de vida da sessão: início, detecção de vitória (morte do
chefe) ou derrota (morte do jogador), e encerramento com exportação automática
de métricas.

% -----------------------------------------------------------------------------
\section{Arquitetura e Ferramentas}
\label{sec:arquitetura-ferramentas}

O desenvolvimento do protótipo apoiou-se em um conjunto de ferramentas
selecionadas por sua adequação ao domínio de jogos 2D com aprendizado por
reforço. A \autoref{tab:ambiente-tecnico} sintetiza as versões utilizadas.

\begin{table}[htb]
    \centering
    \caption{Ambiente técnico utilizado no desenvolvimento e treinamento.}
    \label{tab:ambiente-tecnico}
    \begin{tabular}{l l}
        \toprule
        \textbf{Componente}         & \textbf{Versão / Identificação} \\
        \midrule
        Unity Engine                & 6000.3.17f1                     \\
        Linguagem de scripts        & C\#                             \\
        ML-Agents (pacote Unity)    & 4.0.3                           \\
        Input System (pacote Unity) & 1.19.0                          \\
        Test Framework (Unity)      & 1.6.0                           \\
        Python                      & 3.10.12                         \\
        \texttt{mlagents} (Python)  & 1.1.0                           \\
        PyTorch                     & 2.8.0+cpu                       \\
        ONNX                        & 1.15.0                          \\
        \bottomrule
    \end{tabular}
    \fonte{Elaborado pelo autor.}
\end{table}

\subsection{Unity como Motor de Jogo e Ambiente de Simulação}
\label{subsec:unity-motor}

O Unity \cite{unity2024} foi escolhido como motor de jogo por reunir
capacidades de desenvolvimento 2D nativo --- incluindo sistemas de física
(\texttt{Rigidbody2D}, \texttt{BoxCollider2D}), renderização por
\textit{sprites} e \textit{Input System} moderno --- com suporte integrado ao
\textit{toolkit} de aprendizado de máquina ML-Agents
\cite{juliani2020unity}. Essa integração nativa elimina a necessidade de
adaptadores externos entre o ambiente de simulação e o processo de
treinamento, simplificando o pipeline.

Toda a lógica do jogo foi implementada em C\#, organizada em módulos com
responsabilidades claras: \textit{Player}, \textit{Boss FSM}, \textit{Boss
PPO}, Proveniência, Arena e Métricas. Cada módulo comunica-se por meio de
eventos (\textit{delegates} e \texttt{Action<T>} do C\#), seguindo o
princípio de comunicação dirigida por eventos para minimizar o acoplamento
entre componentes.

\subsection{Unity ML-Agents como Ponte para RL}
\label{subsec:ml-agents-ponte}

O Unity ML-Agents \textit{Toolkit} \cite{juliani2020unity} atua como
intermediário entre o ambiente de jogo e os algoritmos de aprendizado por
reforço. No lado Unity, o agente é implementado como um componente C\# que
herda da classe \texttt{Agent}, expondo métodos para coleta de observações
(\texttt{CollectObservations}), execução de ações
(\texttt{OnActionReceived}) e atribuição de recompensas
(\texttt{AddReward}). No lado Python, o treinador executa o algoritmo PPO
\cite{schulman2017ppo} utilizando PyTorch como \textit{backend} de
computação, comunicando-se com o Unity via protocolo gRPC.

Ao término do treinamento, o modelo é exportado em formato ONNX
\cite{onnx2019}, permitindo inferência diretamente no Unity sem dependência
do processo Python. Essa separação entre fase de treinamento e fase de
inferência é essencial para a avaliação: o modelo avaliado opera de forma
autônoma, determinística e reprodutível.

% -----------------------------------------------------------------------------
\section{Pipeline de Treinamento e Coleta}
\label{sec:pipeline}

O pipeline de desenvolvimento segue um fluxo sequencial de seis etapas, cada
uma produzindo artefatos consumidos pela etapa seguinte. A
\autoref{fig:pipeline} ilustra esse fluxo de forma esquemática.

 \begin{figure}[htb]
     \centering
     \includegraphics[width=0.85\textwidth]{imagens/pipeline-treinamento}
     \caption{Pipeline de treinamento e coleta: do \textit{gameplay} humano
              à avaliação quantitativa FSM vs.\ PPO.}
     \label{fig:pipeline}
     \fonte{Elaborado pelo autor.}
 \end{figure}

\subsection{Etapa 1 --- Coleta de Sessões contra o Chefe FSM}
\label{subsec:etapa1}

Na primeira etapa, o desenvolvedor joga múltiplas sessões como Elian contra o
chefe controlado por FSM. O chefe FSM opera com quatro estados determinísticos
--- \textit{Idle}, \textit{Chase}, \textit{Attack} e \textit{Retreat} ---
com parâmetros fixos de velocidade, alcance de ataque e \textit{cooldown}. A
previsibilidade do FSM é uma característica desejável nesta etapa, pois
garante que as sessões coletadas reflitam o comportamento do jogador humano
em condições estáveis e reprodutíveis.

\subsection{Etapa 2 --- Registro de Proveniência}
\label{subsec:etapa2}

Durante cada sessão, o sistema de proveniência registra automaticamente os
eventos de \textit{gameplay} como nós de um grafo causal. Cada nó contém
atributos padronizados: identificador do evento (\texttt{eventId}), marcação
temporal (\texttt{timestamp}), identificador do ator (\texttt{actorId}), tipo
de ação (\texttt{actionType}), posição no espaço (\texttt{position}), valor
associado (\texttt{value}) e referência ao evento causador
(\texttt{parentEventId}). Essa última referência é a ligação causal que
transforma a sequência temporal de eventos em uma cadeia causal explícita.

O componente \texttt{ProvenanceLogger} escuta eventos emitidos pelos scripts
de \textit{gameplay} --- \texttt{PlayerController}, \texttt{PlayerCombat},
\texttt{PlayerHealth}, \texttt{BossHealth} e \texttt{BossFsmController} ---
e registra cada um no grafo em memória, gerenciado pelo componente
\texttt{ProvenanceGraph}. Ao término da sessão, o
\texttt{ProvenanceExporter} serializa o grafo completo em formato JSON.

\subsection{Etapa 3 --- Extração de Sequências Vitoriosas}
\label{subsec:etapa3}

A extração de padrões comportamentais a partir dos grafos de proveniência é
realizada por um script Python dedicado. Esse script opera em três passos
sequenciais:

\begin{enumerate}
    \item \textbf{Filtragem}: lê os arquivos JSON exportados e seleciona
          apenas as sessões cujo campo \texttt{result} é igual a
          \texttt{"victory"}, descartando sessões encerradas com derrota ou
          incompletas;
    \item \textbf{Localização de eventos-alvo}: percorre os nós de cada
          grafo filtrado e identifica todas as ocorrências do evento
          \texttt{BossDamageTaken}, que indica que o chefe recebeu dano;
    \item \textbf{Extração de sequências}: para cada evento
          \texttt{BossDamageTaken}, retrocede na cadeia causal e coleta as
          três ações imediatamente anteriores pertencentes ao vocabulário de
          ações permitidas (\texttt{PlayerJump}, \texttt{PlayerAttack},
          \texttt{PlayerDash}), formando uma sequência de tamanho fixo.
\end{enumerate}

A \autoref{tab:config-extracao} resume os parâmetros utilizados nesse
processo.

\begin{table}[htb]
    \centering
    \caption{Configuração do processo de extração de sequências vitoriosas.}
    \label{tab:config-extracao}
    \begin{tabular}{l l}
        \toprule
        \textbf{Parâmetro}    & \textbf{Valor}                                  \\
        \midrule
        Ações permitidas      & \texttt{PlayerJump}, \texttt{PlayerAttack},
                                \texttt{PlayerDash}                              \\
        Tamanho da sequência  & 3                                                \\
        Evento-alvo           & \texttt{BossDamageTaken}                         \\
        Formato de saída      & JSON (\texttt{winning\_sequences.json})          \\
        \bottomrule
    \end{tabular}
    \fonte{Elaborado pelo autor.}
\end{table}

O resultado é um arquivo JSON contendo a lista de sequências de ações que, em
sessões vitoriosas, precederam causalmente eventos de dano no chefe. Esse
artefato não constitui uma demonstração direta de \textit{gameplay}: ele
codifica \textit{padrões} de ação associados a desfechos favoráveis, sem
registrar trajetórias completas de estados e ações.

\subsection{Etapa 4 --- Treinamento PPO com Reward Shaping por Proveniência}
\label{subsec:etapa4}

O agente PPO foi treinado utilizando o Unity ML-Agents com o algoritmo PPO
como treinador principal \cite{schulman2017ppo}. A configuração do
treinamento foi definida em um arquivo YAML consumido pelo treinador Python.
A \autoref{tab:config-ppo} apresenta os hiperparâmetros utilizados.

\begin{table}[htb]
    \centering
    \caption{Hiperparâmetros do treinamento PPO.}
    \label{tab:config-ppo}
    \begin{tabular}{l r}
        \toprule
        \textbf{Hiperparâmetro}          & \textbf{Valor}  \\
        \midrule
        \texttt{trainer\_type}           & PPO             \\
        \texttt{batch\_size}             & 1\,024          \\
        \texttt{buffer\_size}            & 10\,240         \\
        \texttt{learning\_rate}          & $3{,}0 \times 10^{-4}$ \\
        \texttt{beta} (entropia)         & $5{,}0 \times 10^{-3}$ \\
        \texttt{epsilon} (clipping)      & 0,2             \\
        \texttt{lambd} (GAE-$\lambda$)   & 0,95            \\
        \texttt{num\_epoch}              & 3               \\
        \texttt{hidden\_units}           & 128             \\
        \texttt{num\_layers}             & 2               \\
        \texttt{max\_steps}              & 500\,000        \\
        \texttt{time\_horizon}           & 64              \\
        Sinal de recompensa              & Extrínseco      \\
        \bottomrule
    \end{tabular}
    \fonte{Elaborado pelo autor.}
\end{table}

Durante o treinamento, o componente \texttt{ProvenanceRewardShaper} mantém um
\textit{buffer} circular das últimas ações executadas pelo agente. A cada
passo de decisão, o \textit{buffer} é comparado com as sequências vitoriosas
carregadas do arquivo JSON. Quando há correspondência exata entre o
\textit{buffer} e alguma sequência vitoriosa, uma recompensa adicional de
$+2{,}0$ é concedida ao agente via \texttt{AddReward}. Essa recompensa
auxiliar opera em conjunto com a recompensa extrínseca do ambiente, sem
substituí-la.

Um aspecto de projeto relevante é a correspondência de vocabulário entre as
ações do agente PPO e as ações humanas registradas na proveniência. As ações
discretas do agente --- pulo, ataque e \textit{dash} --- registram no
\textit{Reward Shaper} as mesmas \textit{strings} utilizadas pelo sistema de
proveniência (\texttt{PlayerJump}, \texttt{PlayerAttack},
\texttt{PlayerDash}), garantindo compatibilidade direta na comparação de
sequências.

O \textit{run} final de treinamento, identificado como
\texttt{boss\_v1\_500k}, alcançou 500.000 passos de treinamento e produziu o
modelo \texttt{BossAgent.onnx}, que foi exportado e integrado ao projeto
Unity em modo \texttt{InferenceOnly} para a etapa de avaliação.

\subsection{Etapa 5 --- Integração e Preparação para Avaliação}
\label{subsec:etapa5}

O modelo ONNX resultante foi integrado ao projeto Unity e associado ao
\textit{prefab} do chefe PPO. Antes da coleta oficial de dados, foi
necessário ajustar as cenas de avaliação para garantir condições equivalentes
entre as duas condições experimentais:

\begin{itemize}
    \item Adição do componente \texttt{PlayerHealth} com 100 pontos de vida
          ao personagem Elian, permitindo que sessões terminassem por derrota
          do jogador;
    \item Implementação do componente \texttt{BossAttackDamage}, que aplica
          dano real de 10 pontos quando o chefe (FSM ou PPO) ataca dentro do
          alcance, com \textit{cooldown} de 0,8 segundos;
    \item Equalização de parâmetros de dano e alcance entre as cenas FSM e
          PPO, eliminando vieses de configuração.
\end{itemize}

\subsection{Etapa 6 --- Avaliação Quantitativa}
\label{subsec:etapa6}

A etapa final consiste na coleta controlada de sessões jogáveis para
comparação. O protocolo de avaliação define a realização de 10 sessões contra
o chefe FSM e 10 sessões contra o chefe PPO, todas jogadas pelo
desenvolvedor em condições padronizadas. As métricas são exportadas
automaticamente em formato CSV pelo sistema \texttt{GameMetrics} e
\texttt{GameMetricsExporter}, incluindo duração da sessão, dano recebido
pelo chefe, dano recebido pelo jogador, contagem de ações do jogador e
contagem de ações do chefe.

A análise dos dados é apresentada nos capítulos de Avaliação e Resultados.

% -----------------------------------------------------------------------------
\section{Protocolo de Avaliação}
\label{sec:protocolo-avaliacao}

A avaliação quantitativa segue um protocolo controlado, projetado para
maximizar a comparabilidade entre as duas condições experimentais dentro das
restrições de escopo do trabalho.

\subsection{Condições Experimentais}
\label{subsec:condicoes}

Duas condições foram definidas:

\begin{itemize}
    \item \textbf{Condição FSM}: o jogador enfrenta o chefe controlado por
          Máquina de Estados Finitos, na cena \texttt{ArenaChefe.unity};
    \item \textbf{Condição PPO}: o jogador enfrenta o chefe controlado pelo
          modelo PPO treinado, em modo \texttt{InferenceOnly}, na cena
          \texttt{BossArena\_PPO.unity}.
\end{itemize}

Em ambas as condições, o personagem Elian é controlado pelo desenvolvedor com
as mesmas mecânicas e parâmetros. Os parâmetros de dano do chefe (10 pontos
por ataque, alcance de 1,25 unidades, \textit{cooldown} de 0,8 segundos)
e a vida do jogador (100 pontos) são idênticos nas duas cenas, garantindo
que diferenças observadas nas métricas reflitam exclusivamente o
comportamento do chefe.

\subsection{Tamanho da Amostra}
\label{subsec:amostra}

A amostra oficial é composta por 20 sessões: 10 na condição FSM e 10 na
condição PPO. Embora o tamanho da amostra seja limitado, ele é suficiente para
identificar diferenças comportamentais mensuráveis entre as duas condições,
que é o objetivo do trabalho. Não se pretende, com essa amostra, estabelecer
significância estatística robusta ou generalizar resultados para outros
contextos --- essas restrições são discutidas como limitações no
\autoref{cap:discussao}.

\subsection{Métricas Coletadas}
\label{subsec:metricas}

As métricas são registradas automaticamente pelo sistema
\texttt{GameMetrics} e exportadas em CSV com formato invariante. As métricas
principais são:

\begin{itemize}
    \item \textbf{Resultado da sessão}: vitória, derrota ou encerramento
          prematuro;
    \item \textbf{Duração da sessão} (em segundos);
    \item \textbf{Dano total recebido pelo chefe}
          (\texttt{boss\_damage\_taken});
    \item \textbf{Dano total recebido pelo jogador}
          (\texttt{player\_damage\_taken});
    \item \textbf{Contagem de ações do jogador}: pulos, ataques e
          \textit{dashes};
    \item \textbf{Contagem de ações do chefe}: ataques realizados; para o
          PPO, inclui-se também a distribuição completa das seis ações
          discretas (Idle, Move Left, Move Right, Jump, Attack, Dash).
\end{itemize}

% -----------------------------------------------------------------------------
\section{Decisões de Escopo e Limitações Metodológicas}
\label{sec:decisoes-escopo}

O escopo do trabalho foi delimitado por decisões deliberadas, tomadas ao
longo do desenvolvimento para manter a viabilidade do TCC sem comprometer a
integridade dos resultados. Esta seção documenta as decisões mais relevantes
e suas justificativas.

\subsection{Remoção da Aprendizagem por Imitação (GAIL)}
\label{subsec:remocao-gail}

O plano inicial do trabalho previa a utilização de \textit{Generative
Adversarial Imitation Learning} (GAIL) \cite{ho2016gail} como mecanismo
complementar de treinamento, a partir de demonstrações gravadas em formato
\texttt{.demo} do ML-Agents. Essa abordagem foi removida do escopo por duas
razões:

\begin{enumerate}
    \item O GAIL introduziria uma dependência de demonstrações diretas,
          conflitando com a proposta central de utilizar proveniência --- e
          não imitação --- como fonte de orientação para o aprendizado;
    \item A coexistência de GAIL e \textit{Reward Shaping} por proveniência
          dificultaria a atribuição de efeitos: seria impossível determinar
          se o comportamento aprendido resultou da imitação, do
          \textit{Reward Shaping} ou de sua combinação.
\end{enumerate}

A remoção do GAIL reforça a contribuição acadêmica do trabalho ao isolar o
efeito da proveniência como único mecanismo auxiliar de orientação.

\subsection{Remoção de Playtesters Externos}
\label{subsec:remocao-playtesters}

O plano original previa a participação de jogadores externos
(\textit{playtesters}) e a aplicação de questionários de experiência. Essa
etapa foi removida por duas razões de escopo:

\begin{enumerate}
    \item A logística de recrutamento, consentimento e coleta presencial de
          dados ultrapassava o prazo disponível para o TCC;
    \item A inclusão de variáveis subjetivas (percepção de dificuldade,
          satisfação) ampliaria o escopo analítico para além do que o
          trabalho se propõe a endereçar.
\end{enumerate}

Em substituição, a avaliação foi redesenhada como uma coleta quantitativa
controlada, conduzida pelo próprio desenvolvedor. Essa escolha é uma limitação
reconhecida --- o avaliador conhece os padrões de ambos os chefes --- mas
garante padronização absoluta das condições de teste e viabilidade temporal.

\subsection{Modelo Congelado na Avaliação}
\label{subsec:modelo-congelado}

O modelo PPO utilizado na avaliação é o resultado do \textit{run} final de
treinamento (\texttt{boss\_v1\_500k}), exportado como arquivo ONNX e
integrado ao Unity em modo \texttt{InferenceOnly}. Após a exportação, o
modelo não foi retreinado, ajustado ou modificado em nenhuma circunstância,
inclusive após a adição do sistema de vida real do jogador
(\texttt{PlayerHealth}).

Essa decisão foi deliberada: a avaliação deve medir o comportamento do modelo
congelado, não de um modelo ajustado \textit{post hoc} com base nos
resultados observados. Como consequência, o modelo avaliado não observa a
vida real do jogador como entrada --- essa observação era um
\textit{placeholder} na arquitetura utilizada durante o treino. Essa limitação
é documentada e discutida no \autoref{cap:discussao}.

\subsection{Preservação do Comportamento Repetitivo}
\label{subsec:preservacao-spam}

Durante as sessões de avaliação, observou-se que o agente PPO apresenta alta
frequência de ações de ataque (\textit{spam}), resultando em comportamento
visualmente repetitivo. Embora uma intervenção no código pudesse limitar essa
frequência --- por exemplo, adicionando um \textit{cooldown} forçado às
ações --- optou-se por preservar o comportamento original do modelo.

A justificativa é metodológica: o comportamento repetitivo é um resultado
experimental do modelo congelado. Corrigi-lo antes da análise
descaracterizaria a avaliação, introduzindo uma modificação \textit{post hoc}
que alteraria os dados coletados sem refletir uma melhoria genuína no
aprendizado do agente.

\subsection{Síntese das Decisões}
\label{subsec:sintese-decisoes}

A \autoref{tab:decisoes-escopo} sintetiza as decisões de escopo e suas
justificativas.

\begin{table}[htb]
    \centering
    \caption{Decisões de escopo e justificativas.}
    \label{tab:decisoes-escopo}
    \begin{tabular}{p{5.5cm} p{8cm}}
        \toprule
        \textbf{Decisão} & \textbf{Justificativa} \\
        \midrule
        Remover GAIL e demonstrações &
            Preservar a proveniência como único sinal auxiliar, isolando
            a contribuição acadêmica. \\
        \addlinespace
        Remover playtesters e questionários &
            Viabilidade temporal e foco em avaliação quantitativa objetiva. \\
        \addlinespace
        Não retreinar o modelo após adicionar vida real &
            Garantir que a avaliação meça o modelo congelado, sem
            ajustes \textit{post hoc}. \\
        \addlinespace
        Preservar o \textit{spam} de ataque do PPO &
            O comportamento é resultado experimental; corrigi-lo
            descaracterizaria a avaliação. \\
        \addlinespace
        Desabilitar JSON de proveniência na cena PPO &
            A métrica principal é o CSV; o JSON da cena FSM já
            evidencia o mecanismo de proveniência. \\
        \bottomrule
    \end{tabular}
    \fonte{Elaborado pelo autor.}
\end{table}

Essas decisões refletem a prioridade de manter a integridade metodológica do
trabalho dentro das restrições práticas de um Trabalho de Conclusão de Curso.
As limitações decorrentes são reconhecidas e discutidas no
\autoref{cap:discussao}, juntamente com sugestões de trabalhos futuros que
poderiam endereçá-las.
